\section{Two exact boundary presentations}\label{sec:boundary}

The contact complex records how pieces meet one another.  A different
question asks how much exact metric information is retained by a presentation
of one piece's polyhedral boundary.  This section records a controlled
sensitivity result; it is not used to prove the atlas theorem.

Use the standard simplex of \cref{sec:model} and divide squared distances by
$2$, so the tetrahedron has unit edge length.  Boundary edges are classified
by their exact incident-plane pair, convex/reflex status, and squared cosine
of the dihedral angle.  Eight exact classes occur.  For each class we sum the
unit-normalised squared edge lengths.

The two presentations differ only in subdivision convention.
\begin{description}[leftmargin=2.7em,style=nextline]
\item[Atomic presentation.]
Retain the edges of the barycentric boundary triangulation and sum all
noncoplanar elementary intervals without collinear merging.
\item[Maximal-reflex presentation.]
Merge connected collinear intervals precisely when the supporting line, the
unordered oriented incident-plane pair, and the exact reflex/cosine class
agree.
\end{description}

\begin{proposition}[Presentation sensitivity]\label[proposition]{prop:boundary}
Across the $82$ atlas rows, the atomic presentation contains $1{,}692$
counted edges and produces $70$ distinct exact profiles.  The maximal-reflex
presentation contains $1{,}506$ edges and produces $72$ profiles.  On the
full-tetrahedron calibration row, the total unit-normalised squared-edge sum
is $3$ in the atomic presentation and $6$ in the maximal presentation.
\end{proposition}
\begin{proof}
The boundary triangles are reconstructed from the exact barycentric chamber
vertices.  Supporting lines, plane normals, reflex signs, squared cosines,
and squared lengths are rational or quadratic data compared without floating
point.  The atomic lane reproduces all $15$ rows of the earlier convex
appendix; the maximal lane applies the declared merging equivalence
independently.  Equality grouping gives the displayed edge and profile
counts.  For the full tetrahedron each original unit edge is split into two
half-edges: the atomic contribution is
$6\cdot2\cdot(1/2)^2=3$, while maximal merging restores the six unit edges.
\end{proof}

Neither presentation is a Dehn invariant.  Dehn data are linear in edge
length and additive under subdivision, whereas a squared-length sum is
nonlinear.  The numbers $3$ and $6$ therefore describe two explicit
presentation conventions rather than an error and a correction.  Classical
scissors-congruence and Dehn-invariant results remain separate
\cite{Sydler1965}.
